\answer{ex:17:1}{$0.2$~W 对 $81$~\textmu W；原始数据流要求每比特不超过 $50$~pJ。}
\answer{ex:17:2}{$c=0.1$ 时为 $5.6$、$8.7$、$9.2$~bit/s，极限 $9.3$~bit/s；$c=0$、$N=1000$ 时为 $65$~bit/s。}
\answer{ex:17:3}{$B=\log_2N+P\log_2P+(1-P)\log_2\frac{1-P}{N-1}=4.87$、$4.37$、$3.29$ 比特/字符，即 $7.3$、$6.6$、$4.9$~bit/s；要达到 $10$~bit/s，每秒需要 $2.3$ 个字符。}
\answer{ex:17:4}{每个分量的误差方差为 $2b/(Nm^2T)+\sigma_\xi^2$；$N\approx90$ 时两项相等；极限为每个窗口每个分量 $\tfrac12\log_2(1+0.25/0.04)\approx1.4$ 比特。}
\answer{ex:17:5}{当 $\eta_bd^2+\eta_db^2<2$，即大约 $\eta<1$ 时这一对收敛：$0.8$ 时收敛，$1.1$ 时出现周期为 $2$ 的稳定振荡，$1.5$ 时为非周期振荡，$2$ 时发散；两个学习者的速度需要拉开。}
\answer{ex:17:6}{阈值 $\approx4.4\sigma$；$102$~kbit/s，$0.10$~mW，而原始数据流需 $205$~mW，小 $2000$ 倍；只有 $5.7\cdot10^{-5}$ 的计数会饱和（多于 $3$ 个脉冲）。}
\answer{ex:17:7}{开放问题。$\text{SNR}\approx0.9$ 时约 $46$~bit/s，独立噪声时约 $330$~bit/s。如果妨碍的是类似信号的噪声，提取的信息会随通道数增加而饱和；如果是对编码的无知，更好的解码器（例如利用脉冲时间的解码器）会使信息增加。}
